\documentclass[12pt]{article}
\usepackage{amsmath, amssymb, amsfonts, geometry}
\geometry{margin=1in}
\setlength{\parskip}{1em}
\setlength{\parindent}{0pt}

\title{Practice Problem Set 4}
\author{ECON 320 -- Introduction to Mathematical Economics}
\date{Fall 2025}

\begin{document}
\maketitle


\section*{Optimization of Multivariable Functions (Chapter 11)}



1. Define the following terms and explain their roles in optimization:
   \begin{enumerate}
     \item Critical point
     \item Gradient vector
     \item Hessian matrix
     \item Definiteness of a matrix
   \end{enumerate}

2. For a twice-differentiable function \(f(x, y)\), explain how the signs of \(f_{xx}\), \(f_{yy}\), and the determinant of the Hessian
   \[
   H = 
   \begin{vmatrix}
   f_{xx} & f_{xy} \\
   f_{yx} & f_{yy}
   \end{vmatrix}
   \]
   determine whether a stationary point is a maximum, minimum, or saddle point.



3. Find and classify all critical points of the function
   \[
   f(x, y) = 4x^2 + y^2 - 4xy + 2x - 3y.
   \]

4. Determine the local extrema (if any) of
   \[
   f(x, y) = 3x^2y - y^3 - 12x^2 + 12y.
   \]
   Compute and interpret the determinant of the Hessian.

5. Let \(f(x, y, z) = 2x^2 + y^2 + z^2 - 2xy - 4xz + 6yz - 3x + y + 2z.\)
   \begin{enumerate}
     \item Find all critical points.
     \item Evaluate the definiteness of the Hessian at those points.
     \item Classify each as a local maximum, minimum, or saddle.
   \end{enumerate}

6. A firm’s production function is \(Q = x^{0.5}y^{0.5}\).
   \begin{enumerate}
     \item Derive the marginal products \(MP_x\) and \(MP_y\).
     \item Show that \(f(x, y)\) exhibits diminishing marginal returns.
     \item Determine the nature of returns to scale.
   \end{enumerate}



7. A monopolist faces a demand function \(p = 100 - 2Q\) and cost function \(C(Q) = 10Q + 0.5Q^2.\)
   \begin{enumerate}
     \item Express profit \(\pi(Q)\) and find the profit-maximizing output.
     \item Verify second-order conditions using \(d^2\pi/dQ^2\).
     \item Extend this to two products: \(\pi(x, y) = 100x + 80y - x^2 - y^2 - xy - 20(x + y)\). 
     Find and classify the stationary point.
   \end{enumerate}

\hrulefill

\section*{Optimization with Equality Constraints (Chapter 12)}



8. Explain the role of the Lagrange multiplier in economic optimization.  
   What does it represent intuitively in problems of:
   \begin{enumerate}
     \item utility maximization subject to a budget constraint, and
     \item cost minimization subject to a production requirement?
   \end{enumerate}

9. State the second-order sufficient conditions for a constrained extremum using the bordered Hessian.  
   Explain why the factor \((-1)^m\) (where \(m\) is the number of constraints) appears in the test.

\newpage

10. Maximize \(f(x, y) = xy\) subject to \(x + y = 10\).
    \begin{enumerate}
      \item Form the Lagrangian and find the stationary point.
      \item Determine whether it is a maximum or minimum using the bordered Hessian.
    \end{enumerate}

11. A firm produces output using two inputs, \(x\) and \(y\), with the production function
    \[
    Q = 4x^{1/2}y^{1/2}.
    \]
    The firm faces the cost constraint \(C = 2x + y = 100.\)
    \begin{enumerate}
      \item Set up the Lagrangian and solve for \(x, y, \lambda.\)
      \item Compute the bordered Hessian and verify whether the solution minimizes cost.
      \item Interpret the Lagrange multiplier economically.
    \end{enumerate}

12. Find the stationary points of
    \[
    f(x, y, z) = x^2 + y^2 + z^2
    \]
    subject to the constraint \(x + y + z = 3.\)
    Determine whether the solution represents a minimum or maximum.

13. Using the method of Lagrange multipliers, find the maximum and minimum values of
    \[
    f(x, y) = x^2 + y^2
    \]
    subject to \(x + 2y = 4.\)

14. Maximize the utility function \(U(x, y) = x^{0.4} y^{0.6}\)  
    subject to the budget constraint \(4x + 2y = 100.\)
    \begin{enumerate}
      \item Derive the first-order conditions.
      \item Find the optimal consumption bundle.
      \item Verify the nature of the extremum using the bordered Hessian.
    \end{enumerate}


\end{document}
